\documentclass[letterpaper]{article}
\usepackage[submission]{aaai2027}
\usepackage[hyphens]{url}
\usepackage{graphicx}
\urlstyle{rm}
\def\UrlFont{\rm}
\usepackage{natbib}
\usepackage{caption}
\frenchspacing

\usepackage{amsmath,amssymb}
\usepackage{booktabs}

\pdfinfo{
/TemplateVersion (2027.1)
}

\setcounter{secnumdepth}{0}

\title{Supplementary Material}
\author{Anonymous Submission}
\affiliations{}

\begin{document}
\maketitle

This document reports method details, the complete baseline comparison, case accounting, and secondary analyses under the same fixed subsets, backbones, and unmodified benchmark interfaces as the main paper.
It expands the evidence map by connecting each contrast to the mechanism and inferential unit it measures.

\section{Method Details and Audit Rationale}

\begin{table}[tb]
\centering
\small
\begin{tabular}{@{}lp{5.0cm}@{}}
\toprule
Symbol & Meaning \\
\midrule
$\mathcal{K}$ & corpus shared by retrieval-based methods \\
$k_{\mathrm{out}}$ & maximum number of emitted diagnoses \\
$B_t$ & shared belief state at stage $t$ \\
$B_{\mathrm{evid}}$ & maximum selected L1 evidence facts supplied to the updater \\
$N$ & size of the updated-score pool before bounded decoding \\
$r$ / $\sim$ & pairwise semantic predicate / its component closure \\
\bottomrule
\end{tabular}
\caption{Notation used consistently in the main paper and supplement.}
\label{tab:notation}
\end{table}

\paragraph{Case-adaptive organization.}
The Level-1 (L1) partition is constrained to use one case-specific comparison axis, labels broader than a single disease, near-exclusive sibling families, a residual family, and optional low-prior families for cannot-miss alternatives.
It serves as a temporary coordinate system specialized to the current syndrome.
If a high-ranked recalled disease has no suitable parent, one non-reducing gap-fill call may enlarge or add families.
Gap-fill supplies a targeted structural repair for parent absence.

\paragraph{Candidate-relative evidence.}
The deployed anti-prototype selector discounts the most tempting vignette prototype and requires a fact to remain discriminative under comparison with live alternatives; it may abstain.
Each consumed fact updates at most one primary support target and one distinct opposing target among L1 families.
Replacing this selector with a salience selector or removing its discrimination hints changes DA Top-1 by $0.04$ (both $p>0.4$), whereas replacing its case-adaptive input axis sharply reduces fact selection and ranking.
This contrast locates the observed gain in the interaction between the case axis and candidate-relative selection.

\paragraph{Concept-aware compression.}
Synonym crowding causes two distinct failures.
\emph{Slot dilution} pushes competing concepts below the output cutoff; \emph{decision splitting} distributes support across several strings for one concept when the scoring contract expects a canonical name.
Concept-aware compression addresses both by preserving clinical equivalence while reallocating slots to distinct hypotheses.

\paragraph{Equivalence construction.}
The implementation first applies a gold-blind pairwise predicate $r$ to candidate labels at the benchmark's scored granularity.
The predicate is symmetrized, candidates become vertices in an undirected graph, and connected components define the transitive closure $\sim$ used for competition.
Within each component, the highest pre-compression rank is retained as representative; ties preserve the frozen input order.
This is single-linkage closure: two labels may share a component through an intermediate chain even when $r$ does not directly relate the endpoints.
The quotient guarantee applies to the closed relation $\sim$; the pairwise predicate, closure, and representative rule are independently auditable implementation components.

\paragraph{Belief-state transition.}
The persisted state is $B_t=(G_t,q_t,s_t,E_t)$: hierarchy, concept-class assignment, relative leaf scores, and evidence annotations.
Local annotation returns updated leaf fields $B_{t+1}=U(B_t,e_t)$.
The arm without score write-back constructs the global pool from the stale $s_t$ fields; the evidence-conditioned score write-back arm replaces them with the corresponding fields in $s_{t+1}$ before any pool or decoder is run.
These updated scores act as relative ranking quantities rather than calibrated probabilities.

\paragraph{Bounded multi-candidate decoding.}
If family $g$ submits $b_g$ leaves, fixing $b_g{=}1$ makes every local ordering error irreversible and hides the eliminated leaf from global arbitration.
APHHM instead permits $b_g{=}2$ when the local margin is small or evidence is thin, or discards family quotas and decodes from the Top-$N$ updated-score pool; OX uses $N{=}15$ before emitting $k_{\mathrm{out}}{=}5$ names.
The cap sweep scores $0.657$, $0.651$, and $0.649$ for cap${=}1$, cap${=}6$, and cap${=}999$, respectively, making decoding robust across these settings on the current OX subset.

\begin{table}[tb]
\centering
\small
\begin{tabular}{@{}p{2.15cm}p{4.45cm}@{}}
\toprule
Decoder arm & Operation on the same score-updated tree \\
\midrule
Closed panel & Restrict discussion to the Top-$15$ pool and project the five emitted names back into that pool. \\
Updated-score-pool decoder & Decode from the same Top-$15$ pool without the closed-panel discussion step. \\
Direct updated-score Top-5 & Emit the five highest-scored leaves directly, without discussion. \\
\bottomrule
\end{tabular}
\caption{Operational distinction among the OX decoder arms.}
\label{tab:decoder-contract}
\end{table}

\section{Evaluation Protocol Details}

Each evaluation uses the first 100 cases under a fixed benchmark order; the released artifacts include the corresponding subset manifests.
All retrieval arms share the same indices and \path{meta-llama/llama-3.3-70b-instruct} backbone~\cite{llama33}; MCR reasoning recall uses the \texttt{gemini-2.5-flash} endpoint~\cite{gemini25flash} with a fixed prompt and no lexical substitution.
Each point estimate is one fixed evaluation pass per configuration.
Candidate enumeration uses temperature zero, although provider-side nondeterminism may remain; paired resampling is used for uncertainty rather than for averaging repeated system runs.
Gold labels are excluded from generation, evidence selection, state update, and arbitration, and every reported result uses the original benchmark scoring interface.

\begin{table}[tb]
\centering
\small
\begin{tabular}{@{}lrrr@{}}
\toprule
Dataset & Raw & Main scored & Audit $n$ \\
\midrule
DA & 100 & 100 & 20 \\
MCR & 100 & 98 & 42 \\
OX & 100 & 100 & 100 \\
\bottomrule
\end{tabular}
\caption{Case accounting. Audit $n$ denotes the stage audit on DA, perturbable-partition audit on MCR, and state factorial on OX. Two MCR cases lack a common fixed-judge score; endpoint tables state further conditioning explicitly.}
\label{tab:case-accounting}
\end{table}

\section{Complete Baseline Table}

\begin{table*}[t]
\centering
\small
\begin{tabular}{@{}lcccccc@{}}
\toprule
& \multicolumn{3}{c}{DA} & \multicolumn{2}{c}{MCR} & OX \\
Method & @1 & @2 & MRR@2 & Acc@1 & R-Recall & micro-F1 \\
\midrule
\textbf{APHHM} & \textbf{0.71} & \textbf{0.78} & \textbf{0.748} & \textbf{0.50} & \textbf{0.753} & \textbf{0.651} \\
MEDDxAgent & 0.62 & 0.71 & 0.665 & 0.24 & 0.412 & 0.491 \\
MAC & 0.61 & 0.67 & 0.640 & 0.23 & 0.527 & 0.570 \\
Dual-Inf & 0.60 & 0.70 & 0.650 & 0.17 & 0.444 & 0.507 \\
i-MedRAG & 0.60 & 0.67 & 0.635 & 0.19 & 0.482 & 0.419 \\
MDAgents & 0.58 & 0.67 & 0.625 & 0.20 & 0.570 & 0.543 \\
Self-refine & 0.57 & 0.62 & 0.595 & 0.21 & 0.447 & 0.530 \\
Flat rerank & 0.56 & 0.63 & 0.595 & 0.17 & 0.404 & 0.495 \\
CoT+RAG & 0.55 & 0.63 & 0.590 & 0.22 & 0.478 & 0.466 \\
Direct CoT & 0.54 & 0.61 & 0.575 & 0.18 & 0.510 & 0.543 \\
Chain-of-Diagnosis + shared corpus & 0.54 & 0.55 & 0.545 & 0.14 & 0.430 & 0.475 \\
Self-consistent CoT (5 samples) & 0.52 & 0.61 & 0.565 & 0.19 & 0.369 & 0.539 \\
Flat beam search & 0.52 & 0.60 & 0.560 & 0.21 & 0.447 & 0.510 \\
Medprompt-style & 0.52 & 0.57 & 0.545 & 0.18 & 0.294 & 0.522 \\
MedRAG & 0.48 & 0.52 & 0.500 & 0.17 & 0.557 & 0.485 \\
DiagnosisGPT-6B$^{\dagger}$ & 0.14 & 0.14 & 0.140 & --- & --- & --- \\
\midrule
Flat rerank $\times 10$ (RRF) & 0.47 & 0.59 & 0.530 & 0.15 & 0.383 & 0.487 \\
\bottomrule
\end{tabular}
\caption{Every completed baseline arm under the unmodified benchmark interfaces. $^{\dagger}$DiagnosisGPT-6B is shown as a local-backbone reference point.}
\label{tab:all-baselines}
\end{table*}

\section{Call-Count Audit}

The approximately nine-call structural proxy reproduces a tree-derived schedule and provides an intermediate inference-volume point.
The ten-trajectory flat reranking control, denoted Flat rerank $\times 10$ (RRF), fuses ten independently generated retrieve--rerank lists with reciprocal rank fusion.
The resource axis is mean model calls per case, measured from LLM cache entries.

\begin{table}[tb]
\centering
\small
\begin{tabular}{@{}lccccc@{}}
\toprule
Dataset & APHHM & Native & Proxy & Flat $\times 10$ & Ratio \\
\midrule
DA & 94.3 & 2.0 & 9.24 & 92.4 & 0.98 \\
OX & 68.6 & 2.0 & 8.98 & 89.8 & 1.31 \\
MCR & 81.2 & 2.0 & 9.32 & 93.2 & 1.15 \\
\bottomrule
\end{tabular}
\caption{Mean model calls per case ($n{=}100$). Flat rerank $\times 10$ (RRF) is matched within $2\%$ on DA and is call-richer than APHHM on OX and MCR.}
\label{tab:compute-audit}
\end{table}

The ten-trajectory control supplies $0.98\times$, $1.31\times$, and $1.15\times$ the APHHM call count on DA, OX, and MCR, respectively.
The nine-call proxy anchors the lower end of the inference-volume ladder, while the ten-trajectory arm provides the similar-call-count comparison.
Output-token reconstructions show the same qualitative separation for the proxy arm; formal resource accounting in the main paper uses the audited call-count axis.

\section{Statistical Evidence Map}

The contemporaneous contrast family contains the adaptive-axis, semantic-partition, evidence-conditioned score write-back, and selector comparisons, with Holm step-down used as conservative multiplicity control.
The axis comparison uses an exact paired binary test, the evidence-conditioned score write-back comparison uses a paired case-bootstrap that recomputes micro-F1, and the selector comparison serves as a negative control.
The semantic-partition comparison permutes equivalence-class membership under fixed cases; its responsive endpoints were selected through the admissibility audit.
The leaf-injection comparison appears separately as a retrospective diagnostic from the earlier judge configuration.

\begin{table*}[t]
\centering
\small
\begin{tabular}{@{}lllcc l@{}}
\toprule
Contrast & Endpoint & Effect & $p$ & $p_{\mathrm{Holm}}$ & Inference unit \\
\midrule
Adaptive vs.\ random axis pipeline & DA Top-1 & $+0.34$ & $3.1{\times}10^{-7}$ & $1.2{\times}10^{-6}$ & pipeline effect \\
Semantic vs.\ blind merge & MCR any-hit@5 & $+0.097$ & $0.015$ & $0.030$ & partition-level \\
Evidence-conditioned score write-back vs.\ none & OX micro-F1 & $+0.075$ & $4.0{\times}10^{-4}$ & $1.2{\times}10^{-3}$ & case-sampled \\
Case-relative vs.\ salience selector & DA Top-1 & $+0.04$ & $0.48$ & $0.48$ & negative control \\
\midrule
No injection vs.\ leaf injection & DA Top-1 & $+0.30$ & $1.5{\times}10^{-5}$ & --- & historical diagnostic \\
\bottomrule
\end{tabular}
\caption{Reported contrasts and their inferential units. Holm values control the four contemporaneous comparisons; leaf injection is listed as a retrospective diagnostic.}
\label{tab:confirmatory}
\end{table*}

\section{Compression Variants and Structural Bounds}

This section reports variants of the equivalence-compression operator, including structurally invariant variants.
For binary Acc@1, paired binary intervals can be defined from endpoint-discordant cases on $n=98$ commonly scored MCR cases.
Open-MRR@5 is summarized descriptively as a paired nonbinary endpoint.
The analysis plan used a $\pm5$ percentage-point non-inferiority margin for binary accuracy.

The count-matched blind merge and its rank-1-size-matched variant are identities at @1 by construction and are evaluated with the responsive any-hit@5 and open-MRR@5 endpoints; the reported $p=0.015$ samples semantic partitions of the fixed cases.
Calibration-only routing is the one variant whose endpoint discordances run in both directions ($4$ favoring the full model and $2$ favoring the comparison).

\paragraph{Endpoint admissibility audit.}
For each candidate endpoint, we generated $200$ randomized equivalence re-partitions per case and counted cases whose endpoint changed.
Top-1 changed in $0/42$ audited MCR cases and $0/35$ audited DA cases, with null standard deviation exactly zero, as predicted for a deletion operator that retains each class's best-ranked member.
MCR any-hit@5 and open-MRR@5 were responsive and are therefore used for the semantic-compression analysis.
This audit distinguishes responsive endpoints from algebraic identities and defines the target endpoints for independent-subset confirmation.

\paragraph{Semantic structure beyond candidate count.}
Across the execution-site cross, the (candidates per case, accuracy) pairs are $(1.81,0.50)$, $(2.02,0.50)$, $(2.11,0.42)$, and $(4.64,0.50)$, showing that list length alone does not order accuracy.
The count-matched blind merge further preserves the deployed class-count and class-size profile while randomizing membership, yet is worse on both responsive endpoints.
The observed difference therefore tracks clinical equivalence structure beyond list shortening.

\paragraph{Partition-semantics randomization.}
The randomized control changes quotient structure while preserving the aggregate partition profile.
The synonym graph is complete in $0/35$ perturbable DA cases and $0/42$ perturbable MCR cases; $82.8\%$ and $85.1\%$ of draws, respectively, commit at least one quotient violation.
Across draws, the control wrongly merges an average of $1.99$ non-synonymous pairs and splits $1.98$ synonymous pairs.
It leaves $0.79$ synonymous duplicate pairs among surviving representatives, whereas the deployed single-linkage merge guarantees $0.00$.
The resulting permutation $p$-value quantifies sensitivity to class membership across randomized partitions of the fixed cases.

\paragraph{Closed-set binding audit.}
On DA, the compared compression configurations emit identical rank-1 labels in all $89$ gated cases, so the open-set Top-1 difference is exactly zero.
All $12$ option promotions under closed rematching already match the target under open scoring, and promotion is not selective for the target (Fisher $p=0.199$).
The $0.112$ closed-set increase is consequently localized to option binding, while the emitted disease concept remains unchanged.

\paragraph{Concept-identity predicate on DA.}
The concept-identity predicate scores $0.57$ option Top-1 on DA, compared with $0.71$ for the deployed configuration; unconditional synonym merging exceeds it by $0.11$ on the same caliber.
The concept-identity predicate yields more concept clusters than the deployed synonym-oriented predicate in $37$ cases and fewer in $0$.
The deployed advantage is concentrated in $10$ cases where a wider surface-form merge absorbs the target into the rank-1 representative class, which closed-set option rematching then credits.
This cross-caliber dissociation shows how closed-set option binding can reward a wider surface-form class even when clinical ranking is unchanged.

\paragraph{Why any-hit@5 is an identity for most variants.}
The full model delivers fewer than $k_{\mathrm{out}}{=}5$ candidates in $86$ of $98$ cases, so any-hit@5 usually reduces to whether the target is in the surviving set.
Because merging retains the best-ranked member of each class, matched labels remain represented, and seven of nine variants have zero moving cases on this endpoint.
Their identical values are structural identities.

\paragraph{Structural identities.}
Serial merge--calibration composition delivers the identical set, the identical rank-1 leaf, and zero moving cases on all three endpoints in this configuration; removing the L1 soft prior moves one case.
These configurations are therefore invariant or near-invariant on this evaluation subset.

\paragraph{Conditional effect sizes.}
Restricted to cases where the delivered list differs, the full model gains $+0.125$ Acc@1 over unconditional merge ($16$ cases) and $+0.107$ over frequency-matched random routing ($28$ cases).
These conditional estimates describe the changed-output cases alongside the full-sample intervals above.

\begin{table*}[t]
\centering
\small
\begin{tabular}{lccl}
\toprule
Variant & Changed outputs & Acc@1 $95\%$ CI & Structural status \\
\midrule
No decision-time routing & 94 & $[-1.4,+2.0]$ & open-MRR differs descriptively \\
Unconditional merge & 16 & $[-1.4,+2.0]$ & accuracy within margin \\
Calibration only & 82 & $[-2.5,+3.0]$ & bidirectional accuracy changes \\
Serial merge $\rightarrow$ calibrate & 7 & $[0,0]$ & endpoint identity \\
Frequency-matched random routing & 28 & $[-1.3,+3.1]$ & accuracy within margin \\
Count-matched blind merge & 33 & $[0,0]$ & rank-1 identity \\
Blind merge, rank-1 size matched & 26 & $[0,0]$ & rank-1 identity \\
Concept-identity predicate & 37 & $[-1.4,+2.0]$ & accuracy within margin \\
No L1 soft prior & 10 & $[-1.0,+1.0]$ & near-identity \\
\bottomrule
\end{tabular}
\caption{Compression variants against the full model. ``Changed outputs'' counts cases with any delivered-list difference. The interval column reports paired binary Acc@1 bounds and is positive when the full model scores higher; open-MRR@5 is summarized descriptively.}
\label{tab:eq-equiv}
\end{table*}

\section{Axis Cascade and Organization Details}

Empty rankings uniformly reflect selector abstention with zero selected facts on the random axis, localizing them to the terminal state of the selector cascade.

\begin{table}[tb]
\centering
\small
\begin{tabular}{@{}lcccc@{}}
\toprule
Axis & mean $n_{\mathrm{sel}}$ & $n_{\mathrm{sel}}{=}0$ & empty rank & @1 \\
\midrule
Case-adaptive & $4.94$ & $17$ & $1$ & $0.71$ \\
Fixed ICD & $0.65$ & $89$ & $11$ & $0.51$ \\
Random & $0.13$ & $97$ & $45$ & $0.37$ \\
\bottomrule
\end{tabular}
\caption{Selector cascade under three L1 axes on DA ($n{=}100$).}
\label{tab:org-cascade}
\end{table}

\section{State-Consistency Factorial Details}

Table~\ref{tab:state-detail} separates the family-cap variants from the decoder-only comparisons on the same score-updated tree.

\begin{table}[tb]
\centering
\small
\begin{tabular}{@{}p{2.35cm}ccc@{}}
\toprule
Variant & F1 & closed $-$ variant & $95\%$ CI \\
\midrule
Full model (closed) & $0.651$ & --- & --- \\
One candidate per family & $0.657$ & $-0.006$ & --- \\
Uncapped family frontier & $0.649$ & $+0.002$ & --- \\
Direct updated-score Top-$5$ & $0.593$ & $+0.058$ & $[0.028,0.085]$ \\
Updated-score-pool decoder & $0.610$ & $+0.041$ & $[0.013,0.069]$ \\
\bottomrule
\end{tabular}
\caption{OX cap and decoder variants at fixed budget with evidence-conditioned score write-back. Paired case-bootstrap intervals are reported for the decoder alternatives.}
\label{tab:state-detail}
\end{table}

\section{Mechanistic Synthesis}

The audits converge on a single representation-centered account.
The case-adaptive axis creates families that support discriminative evidence selection; equivalence-aware competition converts string-level redundancy into concept-level frontier capacity; and evidence-conditioned score write-back carries local evidence into global arbitration.
Stage attribution then connects each observed error to the state transition or interface responsible for it.
Together, these mechanisms explain why APHHM preserves its end-task margins while using a smaller, semantically denser decision frontier.

\bibliography{references}

\end{document}